% ============================================================
% DRAFT NOTE (results)
% Contains: Per-compressor primary outcomes, ratio-level outcomes, calibration,
% and diagnostic evidence.
% Written now: Table and figure scaffolding only.
% Pending: All results, numeric values, interpretations, and decisions.
% Sources: docs/methodology.md
% ============================================================
\section{Results}

% DRAFT NOTE (primary per-compressor outcomes): Report held-out metrics and
% the decision gate for each compressor without pooling their decisions.
\subsection{Primary per-compressor outcomes}

\begin{table}[t]
\centering
\caption{Held-out primary outcomes. Every value and decision awaits the new
training run.}
\label{tab:main-results}
\begin{tabular}{lccccc}
\toprule
Compressor & Positive damage & Major damage & MSE skill 95\% CI & MAE skill & Decision \\
\midrule
\todo{compressor} & \todo{value} & \todo{value} & \todo{value} & \todo{value} & \todo{value} \\
\todo{compressor} & \todo{value} & \todo{value} & \todo{value} & \todo{value} & \todo{value} \\
\todo{compressor} & \todo{value} & \todo{value} & \todo{value} & \todo{value} & \todo{value} \\
\bottomrule
\end{tabular}
\end{table}

% DRAFT NOTE (per-ratio outcomes): Show the four supported ratios separately
% for each compressor and identify any ratio that determines the decision.
\subsection{Outcomes by removal ratio}

\begin{table}[t]
\centering
\caption{Per-ratio held-out breakdown.}
\label{tab:ratio-results}
\begin{tabular}{lccccc}
\toprule
Compressor & Ratio & MSE skill & MAE skill & Positive-row error & \todo{diagnostic} \\
\midrule
\todo{compressor} & \todo{ratio} & \todo{value} & \todo{value} & \todo{value} & \todo{value} \\
\bottomrule
\end{tabular}
\end{table}

% DRAFT NOTE (calibration): Plot observed against predicted damage bins, with
% prompt-cluster uncertainty, separately for each compressor.
\subsection{Calibration}
\begin{figure}[t]
\centering
\figplaceholder{observed versus predicted damage by predicted-damage bin, one
panel per compressor, with prompt-cluster uncertainty}{results/recovered/ifeval-intervention-v1/}
\caption{\todo{Calibration caption and held-out interpretation.}}
\label{fig:calibration}
\end{figure}

% DRAFT NOTE (diagnostics): Report label prevalence, error slices, prediction
% means, and rank correlation without substituting diagnostics for the gate.
\subsection{Diagnostic slices}
\begin{figure}[t]
\centering
\figplaceholder{label prevalence, error on positive- and non-zero-damage
rows, mean prediction versus mean target, and rank correlation, per
compressor}{results/recovered/ifeval-intervention-v1/}
\caption{\todo{Diagnostic caption and held-out interpretation.}}
\label{fig:diagnostics}
\end{figure}
